%% Master CV - Nicholas Fernie
%% Comprehensive overview of competencies, experience, and achievements.
%% Use as a master reference when tailoring role-specific CVs.
%% Compile with: lualatex main_example.tex

\documentclass[11pt,a4paper,sans]{moderncv}
\moderncvstyle{banking}
\moderncvcolor{blue}

\renewcommand*{\namefont}{\fontsize{34}{36}\bfseries\upshape}
\colorlet{firstnamecolor}{color1}
\colorlet{lastnamecolor}{color1}
\colorlet{namecolor}{color1}
\renewcommand*{\sectionstyle}[1]{{\sectionfont\color{color1}#1}}

\usepackage[utf8]{inputenc}
\AtEndPreamble{\hypersetup{
    colorlinks=true,
    linkcolor=blue,
    filecolor=magenta,
    urlcolor=blue,
    pdftitle={Nicholas Fernie - CV},
    pdfpagemode=UseNone,
}}
\usepackage[scale=0.80]{geometry}
\usepackage{import}
\usepackage{needspace}

% personal data
\name{Nicholas}{Fernie}
\address{Brisbane, Queensland, Australia}{}{}
\email{fernienicholas@gmail.com}
\extrainfo{\href{https://www.linkedin.com/in/nic-fernie-28125084}{LinkedIn}, \href{https://github.com/NFernie}{GitHub}}

\begin{document}

\makecvtitle

\vspace{6pt}
\small{Development geologist with an MPhil in Earth Sciences and more than seven years at Santos, spanning wellsite operations, geomechanics, and development geology on producing assets. Cooper Basin subsurface is a sustained focus across thesis work, a 2024 Basin Research paper, and a public geosteering decision-support tool built on well logs, Petrel surface exports, and horizontal-well target windows.}

\section{Core Competencies}
\vspace{1pt}
\begin{itemize}

\item \textbf{Development geology}: producing-asset geoscience after a wellsite year and a geomechanics year at the same operator.

\item \textbf{Well-log interpretation}: gamma ray, resistivity, cuttings, fluorescence, trajectories, and formation tops in a common TVDSS frame.

\item \textbf{Horizontal wells / geosteering}: decision-support for target windows using Petrel-exported surfaces and analogue Cooper Basin wells.

\item \textbf{Basin context}: Cooper-Eromanga thermal history (MPhil; 2024 Basin Research co-author).

\item \textbf{Structural geology}: first-author Scientific Reports paper on shear-zone reactivation (2018).

\end{itemize}

\section{Professional Experience}
\vspace{3pt}
\begin{itemize}

\item{\cventry{Feb 2021 - Present}{Development Geologist}{Santos Ltd}{Brisbane, Australia}{}{\vspace{1pt}
\begin{itemize}
    \item Development geology on Santos producing assets (title and dates from public profile; confidential well metrics omitted).
    \item Cooper Basin remains the technical through-line: MPhil thermal history, 2024 Cooper-Eromanga paper, and public Stimpee analogue wells used for geosteering.
    \item Integrates well data and geological interpretation to support development and well-planning discussions.
\end{itemize}}}

\vspace{3pt}

\item{\cventry{Feb 2020 - Feb 2021}{Geomechanist}{Santos Ltd}{Brisbane, Australia}{}{\vspace{1pt}
\begin{itemize}
    \item Specialist geomechanics posting supporting well planning.
    \item Bridge year between wellsite operations and development geology.
\end{itemize}}}

\vspace{3pt}

\item{\cventry{Jan 2019 - Jan 2020}{Graduate Wellsite Geologist}{Santos Ltd}{Adelaide, Australia}{}{\vspace{1pt}
\begin{itemize}
    \item Operations geology at the wellsite: geological control during drilling and communication with the operations team.
\end{itemize}}}

\end{itemize}

\section{Independent Projects}
\vspace{1pt}
\begin{itemize}
\item{\cventry{2025 - 2026}{Geosteering Guide}{\href{https://github.com/NFernie/Geosteering_Guide}{github.com/NFernie/Geosteering\_Guide}}{}{}{\vspace{1pt}
Decision-support tool for horizontal wells. Ingests Cooper Basin Stimpee 3/4/6 wells to SQLite; FastAPI + React replay of Stimpee 6 logs against a McKinlay TVDSS target window (45/90 stations in window). Documents a Petrel export workflow (LAS 2.0, ASCII surveys, tops, XYZ point-set surfaces). Independent project, not a Santos production system.
}}
\end{itemize}

\section{Education}
\vspace{1pt}
\begin{itemize}

\item{\cventry{2017-2019}{Master of Philosophy, Geological and Earth Sciences}{University of Adelaide}{Adelaide, Australia}{}{\vspace{1pt}
Thesis: ``Thermal History of Central Australia: Cooper Basin, South Australia \& Anmatjira Range, Northern Territory: Insights from Apatite Fission Track and U-Pb Thermochronology.''
}}

\vspace{3pt}

\item{\cventry{2013-2016}{BSc (Hons) Geology / BSc Earth Sciences}{University of Adelaide}{Adelaide, Australia}{}{\vspace{1pt}
Honours thesis (2016): apatite thermochronology of the Bole-Nangodi Shear Zone, northern Ghana.
}}

\end{itemize}

\section{Languages}
\vspace{1pt}
\begin{itemize}
\item English (native).
\end{itemize}

\section{Publications}
\vspace{1pt}
\begin{itemize}
\item Fernie, N., Glorie, S., Jessell, M., Collins, A. S. (2018). Thermochronological insights into reactivation of a continental shear zone in response to Equatorial Atlantic rifting (northern Ghana). \textit{Scientific Reports}. \href{https://doi.org/10.1038/s41598-018-34769-x}{DOI}
\item Nixon, A., Fernie, N., Glorie, S., et al. (2024). Thermal evolution and sediment provenance of the Cooper-Eromanga Basin: Insights from detrital apatite. \textit{Basin Research}. \href{https://doi.org/10.1111/bre.12843}{DOI}
\item Nixon, A., Glorie, S., Fernie, N., et al. (2022). Intracontinental Fault Reactivation in High Heat Production Areas of Central Australia. \textit{G-cubed}. \href{https://doi.org/10.1029/2022gc010559}{DOI}
\end{itemize}

\section{References}
\vspace{1pt}
\begin{itemize}
\item{Available upon request.}
\end{itemize}

\end{document}
